Machine learning offers a way to represent feedback policies with shared
parameters and fit them from sampled decisions. In sequential control,
however, the quality of a decision depends on the future trajectory it
induces. An action-regression loss measures how closely a policy matches
its targets; the control objective measures accumulated cost. Connecting
these two objectives requires knowing how fitting errors affect future
cost, particularly when one neural policy must compromise across states.

The Hamilton--Jacobi--Bellman (HJB) equation makes this dependence
explicit. For smooth deterministic control, the value gradient enters
the Hamiltonian and converts the effect of an action on the dynamics
into a future-cost contribution. Minimizing the Hamiltonian compares
actions using both immediate and future cost. A fixed, possibly
suboptimal policy instead satisfies a Hamilton--Jacobi (HJ) evaluation
equation without this minimization. Its value gradient supplies a
policy-improvement signal without requiring the optimal value function.

This evaluation--improvement construction gives a teacher--student
interpretation of policy learning. The evaluated policy acts as a teacher:
its future-cost sensitivity defines local action costs and feasible
targets. A parameterized student fits the resulting supervision.
The central question is what survives this conversion from cost
information to training labels. Under input constraints, different cost
sensitivities can produce the same target while assigning different
penalties to student deviations. HJ evaluation identifies the information
available before fitting; the student's representation determines which
part must be retained.

We study this question with known differentiable dynamics and a common
cost. Exact examples show identical constrained targets with opposite
cost effects of their best shared regression fit. Nonempty improvement
sets can also be jointly unrealizable even when the student class
contains an average-cost improvement. These failures occur with complete
coverage and globally solved fitting.

For a known isotropic quadratic teaching model, the missing information
is a normal component at the constrained target. Under a specified
value-difference oracle, preserving all cost comparisons within a fixed
student output space requires exactly
$\operatorname{rank}(P^\top\Gamma)$ scalar queries, where $P$ spans
student outputs and $\Gamma$ spans normal ambiguity. We construct
feasible queries attaining this bound. A simple hybrid
using the smaller of the normal and student dimensions is already
optimal in generic relative position. Further savings require alignment;
a spectral bound controls perturbation error.

Cost-aware imitation, Hamiltonian fitting and improvement from imperfect
oracles are established in AggreVaTeD, MPC-Net and MAMBA
\cite{sun2017aggrevated,carius2020mpcnet,cheng2020mamba}. Related foundations
include guided policy search, distillation and regularized improvement
\cite{levine2013guided,czarnecki2019distilling,abdolmaleki2018mpo,geist2019regularized},
constrained HJB policy iteration~\cite{kundu2024policy}, viscosity-based
operator learning~\cite{lee2025pideeponet} and proximal policy
distillation~\cite{spigler2025ppd}. Constraint-induced ambiguity is studied
in inverse reinforcement learning~\cite{schlaginhaufen2023identifiability};
the kernel criterion used here belongs to classical linear information
recovery~\cite{pinkus1986recovery,foucart2021recovery}. Our contribution
is the specific information requirement obtained by combining teacher
constraint geometry with student representation.

The evidence separates exact information recovery from policy improvement.
A 480-fit grid and 120 direct-interface checks compare feasible queries
with full information, student-space controls and a simple hybrid.
A constructed shared-actuation case tests positive rank below both
ambient dimensions; random embeddings identify where there is no
advantage over the hybrid. A separate 600-fit loss intervention tests
normal restoration, scalar weighting, shuffled information and
constraint/capacity controls. Its initial primary is inconclusive;
independent confirmation finds lower Bellman regret despite higher
action MSE, with heterogeneous effects and uncertain closed-loop benefit.
Conditional transfer bounds, joint-realizability examples and the full
nonlinear audit and set studies are supplied in the appendices.
